\section{Билет 10. Сложение колебаний одного направления. Биения. Амплитуда и частота биений.}

\begin{definition}
\textbf{Биениями} называются периодические изменения амплитуды результирующего колебания, возникающие при сложении двух гармонических колебаний одинакового направления с близкими, но не равными частотами ($\omega_1 \approx \omega_2$).
\end{definition}

\begin{theorem}[Вывод уравнения биений]
Пусть складываются два колебания с равными амплитудами $A$ и близкими частотами $\omega_1$ и $\omega_2$ ($\Delta\omega = |\omega_2 - \omega_1| \ll \omega$):
\begin{equation}
    x_1 = A \cos(\omega_1 t), \quad x_2 = A \cos(\omega_2 t).
\end{equation}
Результирующее колебание $x = x_1 + x_2$ находим с помощью формулы суммы косинусов:
\begin{equation}
    x = 2A \cos\left(\frac{\omega_2 - \omega_1}{2}t\right) \cos\left(\frac{\omega_2 + \omega_1}{2}t\right).
\end{equation}
Обозначив среднюю частоту $\omega = \frac{\omega_1 + \omega_2}{2}$ и разность частот $\Delta\omega = \omega_2 - \omega_1$, получаем:
\begin{equation}
    x(t) = \left[ 2A \cos\left(\frac{\Delta\omega}{2}t\right) \right] \cos(\omega t). \label{eq:beats}
\end{equation}
Так как $\Delta\omega \ll \omega$, множитель в квадратных скобках меняется значительно медленнее, чем $\cos(\omega t)$. Его можно считать \textbf{медленно меняющейся амплитудой}:
\begin{equation}
    a(t) = 2A \left| \cos\left(\frac{\Delta\omega}{2}t\right) \right|. \label{eq:beat_amplitude}
\end{equation}
\end{theorem}

\begin{property}
Из уравнения \eqref{eq:beat_amplitude} следуют основные характеристики биений:
\begin{enumerate}
    \item \textbf{Амплитуда биений} периодически меняется от $0$ до $2A$.
    \item \textbf{Частота биений} $\Omega$ равна разности частот складываемых колебаний:
    \begin{equation}
        \Omega = |\omega_2 - \omega_1| = \Delta\omega. \label{eq:beat_freq}
    \end{equation}
    \item \textbf{Период биений} $T_{\text{б}}$ определяется как промежуток времени между двумя последовательными максимумами амплитуды:
    \begin{equation}
        T_{\text{б}} = \frac{2\pi}{\Omega} = \frac{2\pi}{|\omega_2 - \omega_1|}. \label{eq:beat_period}
    \end{equation}
\end{enumerate}
\end{property}

\begin{remark}
Если амплитуды складываемых колебаний не равны ($A_1 \neq A_2$), амплитуда биений не обращается в ноль, а изменяется в пределах от $|A_1 - A_2|$ до $A_1 + A_2$ (так называемая «глубина биений»).
\par\noindent\textbf{Практическое применение:} Явление биений широко используется при настройке музыкальных инструментов. Настраивая струну по камертону, музыкант добивается исчезновения биений ($\Delta\omega \to 0$), что означает совпадение частот колебаний струны и эталона.
\end{remark}

\vspace{1cm}


\begin{figure}[htbp]
    \centering
    \begin{tikzpicture}
    \begin{axis}[
        xlabel = {$t$},
        ylabel = {$x(t)$},
        xmin = 0, xmax = 10,
        ymin = -2.5, ymax = 2.5,
        domain = 0:10,
        samples = 400,
        width = 0.9\textwidth,
        height = 0.4\textheight,
        grid = both,
        grid style = {dashed, gray!30},
    ]
    % Огибающая (медленная)
    \addplot [red, dashed, thick] {2*abs(cos(deg(0.5*x)))};
    \addplot [red, dashed, thick] {-2*abs(cos(deg(0.5*x))};

    % Результирующее колебание (биения)
    \addplot [blue, thick] {2*cos(deg(0.5*x))*cos(deg(4*x))};
    \addlegendentry{$x(t)$ — биения}

    \end{axis}
    \end{tikzpicture}
    \caption{График биений: быстрое колебание с медленно меняющейся амплитудой (огибающая показана красным пунктиром)}
    \label{fig:beats_graph}
\end{figure}

